\documentclass[11pt,a4paper,sans]{moderncv}
\moderncvstyle{classic}
\moderncvcolor{blue}
\usepackage[utf8]{inputenc}
\firstname{Zein}
\familyname{ELAJAMY}
\title{Ingénieur Simulation Numérique | Éléments Finis \& Flambage Dynamique | R\&D Industrielle}
\phone{+33 7 54 45 47 42}
\email{zeinfrance4@gmail.com}

\begin{document}
\makecvtitle

\section{Résumé}
Ingénieur R\&D spécialisé en simulation numérique et performance énergétique. Éléments finis (Abaqus/Metafor), CFD (Ansys Fluent), jumeaux numériques et Python. Expérience ArcelorMittal : modèle thermique validé à 95\%, audit énergétique automatisé. PFE : flambage dynamique de tôles minces.

\section{Expériences}
\cventry{Année 1 (2023}{Jumeau Numérique Thermique}{ArcelorMittal R\&D}{}{}{
\begin{itemize}
\item Développement d'un modèle thermique hors-ligne en Python pour simuler le refroidissement industriel.
\item Utilisation de corrélations physiques et intégration de données géométriques.
\end{itemize}
}
\cventry{Avril 2026}{Projet de Fin d'Études : Flambage Dynamique des Tôles Minces Laminées}{ArcelorMittal R\&D}{}{}{
\begin{itemize}
\item Modélisation de cas tests sous Abaqus et Metafor.
\item Réalisation d'études de sensibilité numérique (maillage, type d'éléments, conditions de symétrie).
\end{itemize}
}
\cventry{Année 2 (2024}{Jumeau Numérique Énergétique}{ArcelorMittal R\&D}{}{}{
\begin{itemize}
\item Développement d'un module de calcul en Python pour les bilans énergétiques.
\item Implémentation d'une logique de calcul multi-niveaux et inversion mathématique de corrélations physiques.
\end{itemize}
}
\cventry{Mai}{Stagiaire Recherche - Thermodynamique des Alliages}{Institut Jean Lamour (IJL)}{}{}{
\begin{itemize}
\item Utilisation de logiciels spécialisés (ThermoCalc) pour le calcul de diagrammes de phases.
\item Analyse et interprétation des résultats expérimentaux.
\end{itemize}
}

\end{document}
